\subsection*{Reading two human questions from one image}
We paired human-authored Chinese questions about a single passage and reversed their order. The two inputs contain the same question glyphs, but a correct response must exchange two distinct answers. On held-out passages from the CMRC training release, a learned planner coupled to the original passage reader produced 437 complete answers among 6,086 requests. It answered 217 of 3,043 question pairs correctly in both orders. Supplying the original question boundaries to the earlier composer gave 440 exact answers. Only two targets occurred as a single contiguous span in the passage.

Further training of the passage reader reduced mean training loss from 2.34 to 1.87 and yielded 423 complete answers. It gained 107 requests and lost 124 relative to the supplied-boundary composer. The paired difference was $-0.28$ percentage points (95\% passage-bootstrap interval, $-1.20$ to $0.62$). The planner recovered the entire reference action sequence on 6,077 requests. Canonicalizing the two questions before inference left 212 correct answers and zero pairs correct in both orders. One passage yielded the founding year 1999 and the shareholder Veron International Limited in either requested order. In the first failed request, the reader copied a description of the author instead of the name and an education history instead of the relationship (Fig.~\ref{fig:paired-passage}).

On 836 passages from the official development release, the original passage reader with the learned planner answered both orders correctly for 518 of 4,800 pairs. Shortening the instruction reduced this to zero. At matched numbers of questions and optimizer steps, training on six different instruction prefixes raised both-order accuracy on a further Chinese wording from zero to 461 pairs. It preserved 516 pairs on the familiar instruction, but still gave zero on the shorter instruction. A previously unseen English instruction followed by Chinese questions gave 371 both-correct pairs with the fixed-prefix planner and 369 with the varied-prefix planner. These tests measure transfer of the instruction format. All four formats request the same operation on the same pair of questions.

A physical trace identified the cause of the zero score. With the shorter Chinese instruction, the varied planner isolated the first question on 9,571 of 9,600 requests and the second on 9,573. It also read fragments of the instruction as additional questions. We therefore decoded the same action scores subject to a contiguous-question constraint. The unchanged weights now gave both correct orders on 505 of 4,800 pairs with the shorter wrapper and 517 with the further Chinese wording. The familiar wrapper gave 518 and the English wrapper gave 371. On individual official questions the decoder returned 999 exact images among 3,170. The gain localizes an instruction-reading error; most pairs still received an incorrect answer in at least one order.
